% Part: first-order-logic
% Chapter: proof-systems
% Section: introduction

\documentclass[../../../include/open-logic-section]{subfiles}

\begin{document}

\iftag{FOL}
      {\olfileid[pt]{fol}{prf}{int}}
      {\olfileid[pt]{pl}{prf}{int}}

\olsection{Introdução}

As lógicas comumente têm tanto uma semântica quanto um sistema de
!!{derivation}. A semântica trata de conceitos como verdade,
satisfatibilidade, validade e consequência.  O propósito dos sistemas de
!!{derivation} é fornecer um método puramente sintático de estabelecer
consequência e validade.  Eles são puramente sintáticos no sentido de
que uma !!{derivation} em tal sistema é um objeto sintático finito,
geralmente uma sequência (ou outra disposição finita) de !!{sentence}s
ou !!{formula}s.  Bons sistemas de !!{derivation} têm a propriedade de
que qualquer sequência ou disposição dada de !!{sentence}s ou
!!{formula}s pode ser verificada mecanicamente como ``correta''.

Os sistemas de !!{derivation} mais simples (e historicamente os
primeiros) para a lógica de primeira ordem foram \emph{axiomáticos}.
Uma sequência de !!{formula}s conta como uma !!{derivation} em tal
sistema se cada !!{formula} individual nela é ou está entre um conjunto
fixo de ``axiomas'' ou segue de !!{formula}s anteriores na sequência por
uma entre um número fixo de ``regras de inferência''---e pode ser
verificado mecanicamente se uma !!{formula} é um axioma e se segue
corretamente de outras !!{formula}s por uma das regras de inferência.
Sistemas de !!{derivation} axiomáticos são fáceis de descrever---e
também fáceis de tratar meta-teoricamente---mas as !!{derivation}s neles
são difíceis de ler e entender, e também difíceis de produzir.

Outros sistemas de !!{derivation} foram desenvolvidos com o objetivo de
facilitar a construção de !!{derivation}s ou facilitar a compreensão de
!!{derivation}s uma vez completas.  Exemplos são a dedução natural,
árvores de verdade, também conhecidas como provas por tableaux, e o
cálculo de sequentes.  Alguns sistemas de !!{derivation} são projetados
especialmente com a mecanização em mente; por exemplo, o método da
resolução é fácil de implementar em software (mas suas !!{derivation}s
são essencialmente impossíveis de entender). A maioria desses outros
sistemas de !!{derivation} representa !!{derivation}s como árvores de
!!{formula}s em vez de sequências. Isso facilita ver quais partes de
uma !!{derivation} dependem de quais outras partes.

Assim, para uma dada lógica, como a lógica de primeira ordem, os
diferentes sistemas de !!{derivation} darão diferentes explicações do
que significa uma !!{sentence} ser um \emph{teorema} e o que significa
uma !!{sentence} ser !!{derivable} a partir de outras. Seja como for
feito isso (via !!{derivation}s axiomáticas, deduções naturais,
!!{derivation}s de sequentes, árvores de verdade, refutações por
resolução), queremos que essas relações correspondam às noções
semânticas de validade e consequência. Escrevamos $\Proves !A$ para
``$!A$~é um teorema'' e ``$\Gamma \Proves !A$'' para ``$!A$~é
!!{derivable} de~$\Gamma$.'' Seja como for $\Proves$ definido,
queremos que corresponda a $\Entails$, isto é:
\begin{enumerate}
\item $\Proves !A$ se, e somente se, $\Entails !A$
\item $\Gamma \Proves !A$ se, e somente se, $\Gamma \Entails !A$
\end{enumerate}
A direção ``somente se'' acima é chamada de \emph{correção}. Um
sistema de !!{derivation} é correto se a !!{derivability} garante
consequência (ou validade). Todo sistema de !!{derivation} decente tem
que ser correto; sistemas de !!{derivation} incorretos não são úteis de
modo algum. Afinal, o propósito inteiro de uma !!{derivation} é
fornecer uma garantia sintática de validade ou consequência. Provaremos
a correção para os sistemas de !!{derivation} que apresentamos.

A direção conversa ``se'' também é importante: é chamada de
\emph{completude}. Um sistema de !!{derivation} completo é forte o
suficiente para mostrar que $!A$~é um teorema sempre que $!A$~é válida,
e que $\Gamma \Proves !A$ sempre que $\Gamma \Entails !A$.
A completude é mais difícil de estabelecer, e algumas lógicas não têm
sistemas de !!{derivation} completos. A lógica de primeira ordem tem.
Kurt G\"odel foi o primeiro a provar a completude para um sistema de
!!{derivation} da lógica de primeira ordem em sua dissertação de 1929.

Outro conceito ligado aos sistemas de !!{derivation} é o de
\emph{consistência}. Um conjunto de !!{sentence}s é chamado inconsistente
se qualquer coisa pode ser !!{derive}da dele, e consistente caso
contrário.  A inconsistência é a contrapartida sintática da
insatisfatibilidade: como conjuntos insatisfatíveis, conjuntos
inconsistentes de !!{sentence}s não fazem boas teorias; são defeituosos
de maneira fundamental. Conjuntos consistentes de !!{sentence}s podem
não ser verdadeiros ou úteis, mas pelo menos passam esse limiar mínimo
de utilidade lógica.  Para diferentes sistemas de !!{derivation}, a
definição específica de consistência de conjuntos de !!{sentence}s pode
diferir, mas como~$\Proves$, queremos que a consistência coincida com
sua contrapartida semântica, a satisfatibilidade. Queremos que seja
sempre o caso que $\Gamma$ é consistente se, e somente se, é
satisfatível. Aqui, a direção ``somente se'' corresponde à completude
(consistência garante satisfatibilidade), e a direção ``se''
corresponde à correção (satisfatibilidade garante consistência).  De
fato, para a lógica clássica de primeira ordem, as duas versões de
correção e completude são equivalentes.

\end{document}
